\section{Methods}
% Mirror the preregistration; cite OSF.
%
% 3.1 Materials.
%   - 7 categories $\times$ 5 commitments $\times$ 2 wordings; Table~\ref{tab:items_examples} gives examples (one per category).
%   - Matching: length, override phrases at 0.40; legitimacy built into COR.
%   - Realism audit $\rightarrow$ test-likeness covariate.
%   - Blinding via opaque IDs.
% 3.2 Models. Table~\ref{tab:models}. Why these: family, scale, reasoning on/off, post-training stage.
% 3.3 Procedure. Chain start, statelessness and matched starting constitution; FORCED vs PERMISSIVE; 4 conditions (quote Pellam/Ashgrove briefly); retries and censoring.
% 3.4 Coding. Rubric steps, fate labels, precedence; erosion definition (merges and deletion); judge validation; second judge 25\% subsample.
% 3.5 Hypotheses: H1, H2a/b, H3, H-alt. One line each, as preregistered.
% 3.6 Analysis. cloglog GLMM formula; fallbacks; Holm; TOST bounds set before data. Simulation validation (type-I error $\TypeI$, power $\PowerS$).
% 3.7 Battery. B1 design; AAR/URR; installed constitutions; inclusion rule.
% 3.8 Integrity and deviations. One paragraph. Point to App.~\ref{app:deviations} for resume-bug forensics, dedup, hazard-join bug caught before unblinding, and partial unblinding. Honesty here is a strength.
